\title{Byte Latent Transformer: Patches Scale Better Than Tokens}

\begin{abstract}
We present the Byte Latent Transformer (\textsc{BLT}), a novel large language model architecture operating at the byte level that, for the first time, achieves parity in performance with traditional tokenization-based LLMs at scale, while offering considerable gains in both inference speed and resilience.
\textsc{BLT} processes inputs by encoding sequences of bytes into patches whose sizes adapt dynamically; these patches become the central computational units. The patch segmentation procedure is governed by measuring the entropy of the subsequent byte, thereby directing greater computational resources and model focus to regions of increased data complexity. We provide the inaugural scaling analysis for byte-level models controlled for floating-point operations, investigating up to 8B parameters and training on 4T bytes. Our empirical findings establish the practicality of training scalable models directly on byte streams, eliminating the requirement for predetermined vocabularies. Efficiency in both training and inference benefits from adaptively choosing longer patches when the input data is more predictable; this approach also yields qualitative gains in reasoning tasks and in the handling of rare generalization cases. On the whole, \textsc{BLT} demonstrates superior scaling trends compared to conventional tokenization-based models under equivalent inference resource constraints, accomplished by concurrent expansion in both patch length and model capacity.
\end{abstract}

\section{Introduction}
\label{section:intro}

We present the Byte Latent Transformer~(\textbf{BLT}), an architecture that operates directly on raw byte sequences without the need for tokenization. For the first time, BLT attains the scalability and performance levels of traditional models that rely on tokenization, while also offering marked gains in both computational efficiency and resilience (\S\ref{sec:robustness}).  
Conventional large language models (llms) are almost entirely trained in an end-to-end manner, with the notable exception of the tokenization phase. This preliminary procedure heuristically segments byte streams into a fixed inventory of tokens, a design choice that affects the manner in which input strings are compressed. This introduces several limitations, including a susceptibility to domain or modality biases~\citep{dagangetting}, increased vulnerability to input perturbations~(\S\ref{sec:robustness}), restricted awareness of orthographic features~\citep{edman2024cute}, and unequal representation across languages~\citep{liang2023xlm,petrov2024language,limisiewicz2024myte}.

Historically, tokenization has played a crucial role since training llms directly on raw bytes incurs prohibitive computational expenses due to the extended sequence lengths required~\citep{xue2022byt5}. Earlier research has addressed this challenge by introducing more computationally efficient self-attention mechanisms~\citep{el2020characterbert, clark2022canine} or even adopting architectures that forgo attention altogether~\citep{wang2024mambabyte}~(\S\ref{section:related_work}). Yet, such advancements mainly facilitate the training of \textit{small models}. When dealing with large models, the main computational bottleneck lies in the sizable feed-forward network layers applied to each byte, rather than in the attention modules.

To optimize computational resource distribution, we introduce an adaptive and learnable approach for aggregating bytes into \textit{patches}~(\S\ref{section:patching}), as well as a novel model architecture that integrates both byte-level and patch-level signals. In contrast to conventional tokenization, {\textsc BLT} does not rely on a predefined patch vocabulary. Instead, flexible byte groupings are transformed into latent patch representations using efficient, trainable encoder and decoder components. Our findings demonstrate that this strategy leads to a \textit{greater} efficiency in compute usage compared to models built on traditional tokenization.

Tokenization-based LLMs assign identical computational resources to each token, which often results in a compromise between computational efficiency and model effectiveness. This occurs because the employed compression heuristics used for inducing tokens may not align with the true complexity inherent in prediction tasks. The foundation of our proposed architecture is the principle that computational effort should be adaptively distributed according to specific needs. For instance, predicting the ending of common words generally constitutes a straightforward, low-entropy task that does not necessitate the involvement of a large transformer, especially in comparison to generating the initial word of a new sentence. This concept is operationalized within {\textsc BLT}'s design~(\S\ref{section:architecture}), which features three transformer components: a pair of lightweight byte-level \textit{local models} and one substantial, globally-scoped \textit{latent transformer}~(\autoref{fig:arch_high_level}). 

In order to determine the optimal grouping of bytes into patches for the purpose of judicious computational distribution, {\textsc BLT} segments the input based on the entropy associated with next-byte predictions. This procedure creates groupings of bytes, or patches, that exhibit relatively homogeneous information density within their respective local contexts.

We conduct the first scaling analysis of byte-level models under a flop-controlled regime, scaling up to 8B parameters and processing up to 4T bytes of training data. Our findings demonstrate that large-scale models can be efficiently trained from raw bytes in an end-to-end fashion, bypassing the need for fixed-vocabulary tokenization. Importantly, {\textsc BLT} achieves parity in training flop-controlled performance\footnote{The computational demands of our models are quantified by counting Floating Point OPerations (flops).} with Llama 3, yet requires up to 50\% fewer flops during inference~(\S\ref{section:scaling}).
Furthermore, we provide evidence that modeling at the byte level leads to notable gains in capturing rare or unusual patterns within the data. Compared to models reliant on tokenizers, {\textsc BLT} exhibits stronger robustness to noisy input and demonstrates improved handling at the character level, as showcased in evaluations targeting orthographic, phonological, and low-resource machine translation tasks~(\S\ref{sec:robustness}).
Lastly, our design allows for concurrent expansion of both model size and patch size while holding the inference flop cost constant. Utilizing longer patch sizes generally reduces the frequency at which the global latent transformer is invoked, allowing these computational savings to be redirected towards increasing overall model capacity. Through experiments that control for inference flops~(\autoref{fig:fixed_inference_scaling}), we report significantly more favorable scaling outcomes compared to traditional tokenization-based approaches.

In conclusion, this work offers several key contributions: 1) We present {\textsc BLT}, a byte latent llm framework capable of adaptively assigning computation to enhance flop efficiency; 2) We demonstrate that our approach achieves flop-matched training performance comparable to Llama 3 up to the 8B parameter scale, and further allows for deliberate trade-offs between slight declines in evaluation metrics and up to 50\% improvements in flop efficiency; 3) The {\textsc BLT} approach introduces a novel scaling paradigm for llms, permitting model size increases under constant inference computational constraints; 4) Our findings reveal that {\textsc BLT} models exhibit increased resilience to input noise and possess a finer-grained understanding of sub-word features in the data, aspects often overlooked by token-centric llms. The code for both training and inference with {\textsc BLT} is publicly available at~\url{https://github.com/facebookresearch/blt}.

\section{Patching: From Individual Bytes to Groups of Bytes}
\label{section:patching}

Dividing a byte sequence into \textit{patches} enables {\textsc BLT} to flexibly adjust computational resources according to contextual requirements. In Figure~\ref{fig:patching_types}, several distinct approaches to partitioning bytes into patches are depicted. Mathematically, the patching function $f_p$ transforms a byte sequence $\pmb{x}=\{x_i,|i=1,\ldots n\}$ of length $n$ into a sequence of $m < n$ patches $\pmb{p}=\{p_j|j=1,\ldots,m\}$ by assigning each $x_i$ to either 0 or 1, where a value of 1 signifies the initiation of a new patch. For both token-level and patch-level architectures, the dominant contributor to computational load is the number of Transformer steps, which, in the case of {\textsc BLT}, equates to the total patches generated by the chosen patching function. As a result, the \textit{patch size}—the mean number of bytes contained within a patch—acts as the central metric governing processing costs for both training and inference phases given a particular patching strategy~(\S\ref{section:flops}).

We proceed by presenting three specific patching methodologies: fixed-length byte patching~(\S\ref{section:static-patch}), segmentation at whitespaces~(\S\ref{section:white-patch}), and adaptive patching utilizing entropy measures from a compact byte language model~(\S\ref{section:dyn-patch}). Additionally, we examine the concept of incremental patching and clarify the distinctions between tokenization and patching~(\S\ref{section:bpe-patch}).

\subsection{Strided Patching Every K Bytes}
\label{section:static-patch}

A common and simple approach to organizing bytes is to segment them into uniform patches of size $k$, as exemplified by MegaByte~\citep{yu2023megabyte}. Employing a constant stride is straightforward in both training and inference pipelines, permits direct adjustment of the average patch dimension, and consequently allows for easy regulation of computational cost in terms of flops. Nonetheless, this patching strategy presents notable limitations. Primarily, it fails to allocate computational resources adaptively based on the local informational content: in some cases, a transformer step $j$ might be spent redundantly, such as when predicting large stretches of whitespace in code, or insufficient computation may be applied to regions with high byte density, like mathematical expressions. Additionally, fixed-size patching can result in the inconsistent partitioning of analogous byte sequences—for instance, splitting the same word at different locations—leading to a lack of contextual consistency.

\subsection{Space Patching}
\label{section:white-patch}

\citet{slagle2024spacebyte} presents an uncomplicated yet effective modification to strided patching by introducing new patches at any space-like byte\footnote{
    Space-like bytes refer to any byte excluding latin letters, numeric digits, or utf-8 continuation bytes. Additionally, every patch is required to include at least one non space-like byte. }, which frequently correspond to linguistic unit boundaries across diverse languages. With Space patching, each word is allocated a latent transformer computation step (i.e., greater computational resources), facilitating uniform patching of words within sequences and dedicating additional flops to challenging predictions, which commonly appear after spaces. To illustrate, identifying the initial byte in response to the query ``Who composed the Magic Flute?\uline{\hspace{.6em}}'' poses more difficulty than predicting subsequent bytes after ``M''; this is because the initial character significantly constrains plausible options, thereby making the prediction of the remainder of ``Mozart'' much easier. Nevertheless, space patching exhibits limitations in its adaptability across all languages and domains, and notably cannot dynamically adjust patch size. In the following, we introduce an alternative patching strategy, leveraging the observation that the first bytes of words are generally harder to predict, while simultaneously enabling patch size to be adaptively managed.

\subsection{Entropy Patching: Using Next-Byte Entropies from a Small Byte LM}
\label{section:dyn-patch}

Instead of utilizing rule-based heuristics like whitespace, we employ a data-centric strategy to detect locations where next-byte predictions are highly uncertain. Our method, termed \textit{entropy patching}, determines patch boundaries by leveraging entropy estimates.

We train a small byte-level auto-regressive language model on the training data for {\textsc BLT} and compute next byte entropies under the LM distribution $p_{e}$ over the byte vocabulary $\mathcal{V}$:

\begin{equation}
H(x_i) = \sum_{v \in \mathcal{V}} p_{e}(x_i=v|\pmb{x}_{<i}) \log p_{e}(x_i=v|\pmb{x}_{<i})
\end{equation}

We explore two approaches for detecting patch boundaries based on entropies $H(x_i)$. The first method selects points where entropy exceeds a fixed global threshold, as depicted in~\autoref{fig:patching}. The second approach highlights locations where the entropy is significantly higher compared to the preceding value. This method can also be understood as pinpointing instances where the sequence of entropies within a patch deviates from an approximately monotonic decrease.

\begin{align*}
    \text{Global Constraint}\hspace{1cm}H(x_t)& > \theta_g\\
    \text{Approx. Monotonic Constraint}\hspace{1cm}H(x_t)& - H(x_{t-1}) > \theta_r
\end{align*}

Patch boundaries are determined through a simple preprocessing procedure performed at the dataloading stage. This approach contrasts with the method used by~\citet{nawrot-etal-2023-efficient}, where a classifier is employed to predict patch boundaries based on entropy measurements. In our experiments~(\S\ref{section:experiments}), we evaluate and contrast both techniques for differentiating low entropy from high entropy bytes.

\subsection{The Byte-Pair Encoding (BPE) Tokenizer and Incremental Patching}
\label{section:incremental-patching}
\label{section:bpe-patch}

Contemporary large language models such as our baseline Llama 3 typically employ subword tokenizers, for example bpe~\citep{gage1994new, sennrich-etal-2016-neural}. Throughout this work, we define ``tokens'' as collections of bytes selected from a \textit{finite} vocabulary established before model training, in contrast to ``patches,'' which represent dynamically constructed sequences that are not restricted by a predetermined vocabulary. Importantly, unlike patches, tokens do not provide the model with direct access to the original byte-level information.

An important advancement introduced by {\textsc BLT} in comparison to tokenization-based models is the reconceptualization of the balance between vocabulary size and computational requirements. In traditional llms, enlarging the vocabulary leads to tokens that represent longer text segments on average, which reduces the number of processing steps needed; however, it also necessitates a larger output dimension in the model’s final projection layer. This inherent trade-off constrains tokenization methods, making it challenging for them to substantially modify token size without incurring additional inference overhead. To illustrate, Llama 3 achieves an increase in average token length from 3.7 bytes to 4.4 bytes, but this comes at the expense of expanding its embedding table to four times the size used in Llama 2.

During generation, {\textsc BLT} must determine at each position in the byte sequence whether it has reached a patch boundary, as this governs when the Latent Transformer is applied for additional computation. This determination has to be made without knowledge of the bytes that are yet to be produced in the sequence. Therefore, the process of patching cannot rely on information about subsequent bytes when partitioning the byte sequence. In precise terms, a patching method $f_p$ exhibits the property of incremental patching if it meets the following criterion:
$$f_p(\pmb{x}_{<i}) = f_p(\pmb{x})_{<i}$$
The byte pair encoding (bpe) approach does not fulfill the requirements of incremental patching, since the same sequence prefix may be segmented in varying ways depending on what follows, violating the property above\footnote{Introducing a specific delimiter token to signify patch boundaries can adapt bpe into an incremental patching method, but at the cost of elongating the byte sequence.}.

\section{{\textsc BLT} Architecture}
\label{section:architecture}
The {\textsc BLT} framework consists of a substantial global autoregressive language model which processes patch-level representations. Complementing this, two lightweight local modules are utilized: one for converting byte sequences into patches, and the other for reconstructing byte sequences from patch representations~(\autoref{fig:arch_high_level}).

\subsection{Latent Global Transformer Model}

\textit{The Latent Global Transformer} refers to an autoregressive transformer architecture $\mathcal{G}$ composed of $l_{\mathcal{G}}$ layers, designed to transform a sequence of latent patch representations $p_j$ into corresponding output patch representations $o_j$. In this work, indices $j$ and $i$ are consistently used to identify patches and bytes, respectively. Employing a block-causal attention mask~\citep{dubey2024llama}, the global model enforces attention only up to and including the current patch within each document. Since this component accounts for the majority of the computational cost during both pre-training and inference phases, selective invocation of the model enables dynamic allocation of computational resources across different regions of the input and output, depending on their respective complexities.

\subsection{Local Encoder}
\label{section:local}

\textit{The Local Encoder Model}, represented as $\mathcal{E}$, constitutes an efficient transformer model with significantly fewer layers, $l_{\mathcal{E}} << l_{\mathcal{G}}$, than its global counterpart. Its principal purpose is to rapidly convert a sequence of input bytes $b_i$ into highly expressive patch-level representations, $p_j$. In contrast to conventional transformer designs, this model integrates a cross-attention mechanism following each transformer block, allowing for the aggregation of byte-level representations into patch-level representations (see~\autoref{fig:crossattn}).

Initially, the input byte sequence $b_i$ is transformed into vector embeddings $x_i$ through multiplication with an embedding matrix in $\mathbb{R}^{256 \times h_{\mathcal{E}}}$. Optionally, these embeddings may be further enriched with hash-embeddings, as detailed in~(\S\ref{sec:hash-emb}).

Through a repeated sequence of transformer and cross-attention blocks~(\S\ref{sec:enc-cross-attn}), these embedded representations are mapped into patch representations $p_i$, which subsequently serve as input to the global transformer $\mathcal{G}$. Each transformer block employs a \textit{local block causal} attention scheme, wherein each byte’s attention is confined to a fixed window of $w_{\mathcal{E}}$ preceding bytes. This windowing allows context to span patch boundaries while strictly prohibiting attention across different document boundaries. The forthcoming subsections provide further specification of the embedding process and cross-attention module.

\subsubsection{Encoder Hash n-gram Embeddings}
\label{sec:ngram-emb}
\label{sec:hash-emb}
To produce strong and informative representations at every position $i$, it is important to include contextual data from previous bytes. In {\textsc BLT}, this is addressed by encoding both the specific byte $b_i$ on its own as well as treating it as a component of a byte n-gram. At each step $i$, we begin by generating byte-grams

\begin{align}
    g_{i,n}=\{b_{i-n + 1},\ldots, b_i\}
\end{align}

for every byte position $i$ and for $n$ in the range of three to eight.\footnote{
Byte-grams are excluded for sizes $n$ or greater when $i < n$. }

Subsequently, we define hash $n$-gram embeddings, where every byte $n$-gram is mapped to a position in a dedicated embedding table $E_{n}^{hash}$ of predetermined size using a hash function. This mapping is independently maintained for each $n \in\{3,4,5,6,7,8\}$~\citep{bai2010learning}.  
The resulting $n$-gram embedding is incorporated into the original byte embedding, after which the sum is normalized and fed into the local encoder. The composite embedding is computed as follows:

\begin{align}
e_i &= x_i + \sum_{n = 3,...,8}E_{n}^{hash}(\text{Hash}(g_{i,n})) \\
\text{where, Hash}(g_{i,n}) &= \text{RollPolyHash}(g_{i,n}) \% |E_{n}^{hash}|
\end{align}

To compute $e_i$, we divide by the total number of $n$-gram sizes plus one and utilize the RollPolyHash function as specified in Appendix~\ref{appendix:rollpolyhash}. Section~\ref{section:ablations} presents an analysis of how varying $n$ values and embedding table sizes in $n$-gram hash embeddings influence flop-controlled scaling law behavior. Beyond employing hash-based $n$-gram embeddings, we also conducted experiments with frequency-based $n$-gram embeddings; further information about these experiments can be found in Appendix~\ref{appendix:ngrams}.

\subsubsection{Encoder Multi-Headed Cross-Attention}
\label{sec:enc-cross-attn}
Our approach draws heavily from the cross-attention mechanism used in the Perceiver framework~\citep{jaegle2021perceiver}. However, a key distinction in our design is that the latent variables are associated with representations for variable-sized patches, rather than relying on a predetermined set of latent vectors~(\autoref{fig:crossattn}), and each latent representation only focuses attention on the bytes constituting its specific patch. Each patch $p_j$ is assigned a query vector, derived by applying a pooling operation to the byte representations of $p_j$, and subsequently subjected to a linear transformation defined by $\mathcal{E}_{C} \in \mathbb{R}^{h_{\mathcal{E}} \times (h_{\mathcal{E}} \times U_{\mathcal{E}})}$, where $U_{\mathcal{E}}$ indicates the count of cross-attention heads in the encoder. Specifically, denoting by $f_{\text{bytes}}(p_j)$ the ordered byte sequence for patch $p_j$, the query vectors are computed as follows:

\begin{align}
P_{0,j} &= \mathcal{E}_{C}(f_\text{bytes}((p_j)), f~ \text{is a pooling function} \\
P_l &= P_{l-1} + W_o\left(\text{softmax}\left(\frac{QK^T}{\sqrt{d_k}}\right)V\right) \\
\text{where } Q_j &= W_q(P_{l-1,j}), K_i = W_k(h_{l-1,i}), V_i = W_v(h_{l-1,i}) \\
h_l &= \text{Encoder-Transformer-Layer}_l(h_{l-1})
\end{align}

Here, $P \in \mathbb{R}^{n_p \times h_{\mathcal{G}}}$ denotes the $n_p$ patch representations that serve as inputs to the global model. These representations are constructed by aggregating the byte embeddings $e_i$ associated with each individual patch $p_j$. The matrices $W_q$, $W_k$, $W_v$, and $W_o$ correspond to the linear transformations applied for the queries, keys, values, and outputs, respectively. Notably, the keys and values are derived by projecting the byte representations $h_i$ from the preceding layer (or directly from $e_i$ in the initial layer). A masking mechanism, tailored to patching, restricts each query $Q_j$ to only access keys and values related to the bytes within the same patch $j$. As multi-headed attention operates over $Q, K$, and $V$, and since patch representations often have a higher dimensionality ($h_{\mathcal{G}}$) than $h_{\mathcal{E}}$, $P_l$ is structured as several heads with each having dimension $h_{\mathcal{E}}$ during cross-attention. These head outputs are subsequently concatenated to match the $h_{\mathcal{G}}$ dimensional space. Moreover, queries, keys, and values are all normalized via pre-LayerNorm, and the cross-attention module omits positional embeddings. A residual pathway encompasses the cross-attention layer to facilitate gradient flow.

\subsection{Local Decoder} 

Analogous to the architecture of the local encoder, the local decoder $\mathcal{D}$ employs a compact transformer design, comprising $l_{\mathcal{D}}$ layers such that $l_{\mathcal{D}} << l_{\mathcal{G}}$. Its purpose is to reconstruct raw bytes $y_i$ from the sequence of global patch representations $o_j$. The local decoder generates the sequence of raw bytes conditionally, relying on the sequence of bytes decoded up to the current position, and utilizes the hidden states generated by the local encoder as input. The model processes the data through $l_{\mathcal{D}}$ layers, alternating between cross-attention and standard transformer layers. Each decoder layer first performs cross-attention to map global patch representations to byte-level embeddings, after which a transformer block is applied to these byte representations to model dependencies in the decoded byte sequence.

\subsubsection{Decoder Multi-headed Cross-Attention} 

In the decoder cross-attention, the roles of the queries and key/values are interchanged i.e. the byte-representations are now the queries, and the patch representations are now the key/values. The initial byte-representations for the cross-attention are initialized as the byte embeddings from the last encoder layer i.e. $h_{l_{\mathcal{E}}}$. The subsequent byte-representations for layer $l$, $d_{l,i}$ are computed as:

\begin{align}
D_0 &= h_{l_{\mathcal{E}}} \\
B_l &= D_{l-1} + W_o\left(\text{softmax}\left(\frac{QK^T}{\sqrt{d_k}}\right)V\right), \\
  &\text{where } Q_i = W_q(d_{l-1,i}), K_i = W_k(\mathcal{D}_C(o_j)), V_i = W_v(\mathcal{D}_C(o_j)) \\
  D_l &= \text{Decoder-Transformer-layer}_l(B_l)
\end{align}

In this context, $W_k$ and $W_v$ denote the key and value projection matrices, which carry out linear transformations followed by the split operation $\mathcal{D}_C$ on the global model’s terminal patch outputs, $o_j$. The query projections are realized via $W_q$, mapping from the previous decoder transformer layer’s byte-level representations $d_{l-1}$ (or from $h_{l_{\mathcal{E}}}$ when dealing with the first layer). The projection matrix $W_o$ acts on the resulting output, ensuring that $B$ is of dimension $\mathbb{R}^{h_{\mathcal{D}} \times n_b}$, where $n_b$ indicates the number of generated output bytes. The subsequent decoder layer outputs, $D_l$, arise from running a decoder transformer layer over $B$, the output from the cross-attention procedure. Similar to the method employed for cross-attention in the local encoder, we implement multi-headed attention, utilize pre LayerNorm for normalization, omit positional embeddings, and enclose the cross-attention operation within a residual connection. 

\section{Experimental Setup}
\label{section:experiments}
We construct our experimental design with precision to objectively evaluate {\textsc BLT} alongside models that use tokenization, carefully ensuring that {\textsc BLT} does not benefit from any incidental increase in context window length during comparison.

\subsection{Pre-training Datasets} In this study, we utilize two primary datasets for pre-training models across all examined scales: (1) the Llama 2 dataset~\citep{touvron2023llama}, which contains 2 trillion tokens sourced from a wide array of publicly available texts. These sources are subject to extensive cleaning and filtering processes to enhance data quality; and (2) {\textsc BLT}-1T, a newly compiled corpus comprising 1 trillion tokens. This collection aggregates data from multiple public domains and incorporates a segment of the pre-training dataset released through Datacomp-LM~\citep{li2024datacomp}. The first dataset serves in experiments concerning scaling laws, specifically for identifying the optimal number of tokens according to findings in~\cite{dubey2024llama}, thereby informing the selection of suitable architectures for {\textsc BLT}. The second dataset facilitates a full pre-training cycle intended to benchmark against Llama 3 performance across downstream evaluation tasks. Notably, both datasets exclude any information originating from Meta's products or services. Additionally, for baseline experiments involving tokenizer-based architectures, we employ the Llama 3 tokenizer, which uses a 128K token vocabulary. This tokenizer demonstrated superior baseline outcomes in our assessments compared to the Llama 2 tokenizer.

\subsection{Entropy Model} For our experiments, the entropy model employed is a byte-level language model, trained on the same data distribution as the full {\textsc BLT} model. By default, our configuration utilizes a transformer comprising 100M parameters, spread across 14 layers with a hidden size of 512 and equipped with sliding window attention spanning 512 bytes. All other hyperparameter choices are aligned with those used in our local and global transformers. We evaluated models with varying capacities, receptive field sizes, and architectural changes, as outlined in \autoref{section:ablations}. Notably, when the receptive field is sufficiently limited, it becomes possible to implement the learned entropy model as a compact lookup table.

\subsection{Entropy Threshold and Equalizing Context Length} For entropy-based patching models, we determine a patching threshold to yield a targeted mean \textit{patch size} over the combined pretraining datasets. In {\textsc BLT}, as opposed to traditional tokenization, the selection of \textit{patch size} is flexible, significantly impacting the effective context available to the model. To ensure parity in average context length and prevent potential advantages for models with larger patches, we fix the expected number of bytes per batch. Consequently, models configured with larger patch sizes utilize proportionally shorter sequence lengths. Specifically, on Llama 2 data, we maintain an 8k byte context window, while for the {\textsc BLT}-1T corpus, the average context is increased to 16k bytes, all while keeping the mean batch size at 16M bytes.

Although the overall batch size remains unchanged, dynamic patching approaches result in varying proportions between bytes and patches when assembling data batches. To enhance computational efficiency, our {\textsc BLT} training procedure organizes batches so that each contains an equal number of patches, thereby circumventing the need for padding within the costly latent transformer module. In practice, we standardize byte sequences by padding or truncating them to fixed lengths of 12k and 24k bytes for the Llama 2 and {\textsc BLT}-1T datasets, respectively. This strategy prevents memory overflow caused by atypically large patches within sequences during training.

\subsection{Entropy Model Context}
\label{section:entropy-context}
Our experimental observations indicate that entropy patching tends to produce increasingly larger patches, particularly in structured scenarios such as multiple choice tasks (refer to patching on an MMLU example in \autoref{fig:mmlu_patching}) which frequently display high repetitiveness. These observed differences arise from the reduced entropy associated with repeated segments within the entropy model context. Consequently, for the comprehensive run of {\textsc BLT}-Entropy with a patch size of 4.5, we refresh the entropy context using new lines and implement an approximate monotonicity constraint, as this approach is less susceptible to "entropy drift" resulting from varying context lengths. This adjustment modifies only the entropy calculation process; our approach to determining the entropy threshold remains unchanged.

\subsection{FLOPs Estimation} 
\label{section:flops}

Our approach closely adheres to the transformer flop calculation formulas introduced in Chinchilla~\citep{hoffmann2022training}, accounting for flops associated with the feed-forward networks, qkvo projections in the self-attention module, as well as both the attention computation and the output projection. However, in contrast to their setup, we consider the input embedding layer to operate via a computationally efficient lookup mechanism, rather than a dense matrix multiplication, thus treating it as incurring zero flops. Consistent with earlier studies, we assume the backward pass requires twice as many flops as the forward pass.

To compute flops \textit{per byte} for {\textsc BLT} models, we add up the flops for the local encoder transformer, the global latent transformer, and the local decoder transformer, together with the cross attention blocks in the encoder and the decoder:

\begin{align}
    \text{FL}_{\text{{\textsc BLT}}} &= \text{Transf. FL}(h_{\mathcal{G}}, l_{\mathcal{G}}, m=n_{ctx}/n_p,V=0) / n_p\\
   &+ \text{Transf. FL}(h_{\mathcal{E}}, l_{\mathcal{E}}, m=w_{\mathcal{E}}, V=0) \\
   &+ \text{Transf. FL}(h_{\mathcal{D}}, l_{\mathcal{D}}, m=w_{\mathcal{D}}, V=256) \\ 
    &+ \text{Cross Attn. FL}(h_{\mathcal{E}}, l_{\mathcal{E}}, m=n_p, r=n_p/k) \cdot k/n_p \\ 
   &+ \text{Cross Attn. FL}(h_{\mathcal{D}}, l_{\mathcal{D}}, m=k, r=k/n_p) 
\end{align}

Here, $n_{ctx}$ denotes the number of bytes in the sequence, $n_p$ represents the patch size, $r$ specifies the proportion of query vectors relative to key/value pairs, and $k$ indicates the ratio between the patch dimension and byte dimension—that is, the count of local model divisions concatenated to yield the global model embedding ($k=2$ as depicted in \autoref{fig:crossattn}). The variable $V$ stands for the size of the output projection vocabulary, which is exclusively relevant for the local decoder component. Depending on whether a given module processes the sequence at the byte level or at the patch level, the attention mechanism utilizes a distinct context window, denoted $m$. We adjust our calculation of attention flops to reflect the unique characteristics of each part. Comprehensive formulas for evaluating Transformer-FLOPs and Cross-Attention FLOPs can be found in Appendix~\ref{section:flops-appendix}.

\subsection{Bits-Per-Byte Estimation} Perplexity only makes sense in the context of a fixed tokenizer as it is a measure of the uncertainty for each token. When comparing byte and token-level models, following previous work~\citep{xue2022byt5, yu2023megabyte, wang2024mambabyte}, we instead report Bits-Per-Byte (BPB), a tokenizer independent version of perplexity. Specifically:

\begin{align}
    \text{BPB}(x)&= \frac{\mathcal{L}_{CE}(\pmb{x})}{\ln(2)\cdot n_{\text{bytes}}}
\end{align}

Here, the aggregate cross-entropy loss, which quantifies the model's uncertainty with respect to the dataset $\pmb{x}$, is divided by the total number of bytes in $\pmb{x}$ and by a normalization constant.

\subsection{Transformer Architecture Hyperparameters} 

Across all transformer modules in {\textsc BLT}, encompassing both the local and global models, our design choices are primarily aligned with the specifications in Llama~3~\citep{dubey2024llama}. The feed-forward components utilize the SwiGLU activation function~\citep{swiglu}, and rotary positional embeddings (RoPE)~\citep{RoPe} with $\theta=500000$ as in~\citep{xiong2024effective}, are applied exclusively within self-attention mechanisms. RMSNorm~\citep{RMSNorm} is adopted for layer normalization throughout. For all self-attention layers employing static attention masks, such as \textit{block causal} or \textit{fixed-window block causal}, Flash attention~\citep{dao2022flashattention} is employed, and the attention window is set to 512 for these fixed-width patterns. For cross-attention modules, which require dynamically determined, patch-dependent masking, we leverage Flex Attention\footnote{\url{https://pytorch.org/blog/flexattention}} to generate efficient fused kernels, allowing substantial training speed-up.

\subsection{{\textsc BLT}-Specific Hyperparameters} 

To evaluate the performance of the {\textsc BLT} models, we perform experiments focused on two primary aspects: scaling behavior and results on downstream tasks. For these investigations, we experiment with models of varying parameter sizes: 400M, 1B, 2B, 4B, and 8B, with detailed architectural hyperparameters summarized in Appendix~Table~\ref{tab:arch}.

For the initialization of queries in the initial cross-attention layer within the local encoder, we employ max-pooling. Each {\textsc BLT} model utilizes $500{,}000$ hashes and a single hash function, setting n-gram sizes from 3 to 8. All model variants are trained using a learning rate of $4e-4$. This learning rate selection—identical for both token-based and {\textsc BLT} variants—results from hyperparameter tuning over the $1e-3$ to $1e-4$ range at the 400M and 1B scales, indicating that this configuration yields optimal results.

For experiments exploring scaling trends on Llama-2 data, we adopt the recommended batch sizes given in~\cite{dubey2024llama}, or the equivalent values in bytes. For optimization, we use AdamW~\citep{adamw} with hyperparameters $\beta_1 = 0.9$, $\beta_2 = 0.95$, and $\epsilon = 10^{-8}$. The training schedule includes 2000 steps of linear warm-up, followed by a cosine decay to a final learning rate of zero. In addition, we incorporate a weight decay of 0.1 and perform global gradient clipping with a maximum norm of 1.0.

\section{Scaling Trends}
\label{section:scaling}

We offer a comprehensive analysis of the scaling behaviors characteristic of byte-level models, providing insights to guide further development of {\textsc BLT} architectures. Our scaling investigation is designed to overcome several shortcomings of past studies on byte-level models by: (a) systematically comparing performance under compute-optimal training scenarios, (b) training equivalent 8B-parameter models using substantial datasets—up to 1T tokens or 4T bytes—and evaluating their effectiveness on downstream applications, and (c) quantifying scaling patterns under inference-cost-constrained conditions. In subsequent sections, we will delve into the distinct benefits associated with modeling sequences at the byte level.

\subsection{Parameter Matched Compute Optimal Scaling Trends}
\label{section:parameter-matched-scaling}
Employing the Llama 2 dataset, we conduct training of several \textit{compute-optimal} bpe and {\textsc BLT} models at four distinct parameter scales, spanning from 1B to 8B. We visualize the relationship between training flops and language modeling accuracy using a representative portion of the overall training data blend. The bpe models are optimized according to the best-established ratio between model parameters and dataset size, following the guidance of Llama~3~\citep{dubey2024llama}. Such a \textit{compute-optimal} configuration is expected, in theory, to maximize performance on the given training set under a fixed computation budget~\citep{hoffmann2022training}, establishing a strong baseline for comparison. For every bpe model trained, we construct an equivalent {\textsc BLT} model, ensuring the Latent Transformer employed matches the associated bpe Transformer in both scale and architectural structure, and train both on identical data.

As shown in~\autoref{fig:scaling-compute-opt} (right), {\textsc BLT} models consistently achieve performance on par with, or superior to, their bpe-based equivalents, and this pattern persists as both model size and computational resources are increased. To our knowledge, {\textsc BLT} represents the first instance of a byte-level Transformer architecture exhibiting scaling trends comparable to BPE-based models when operating within compute-optimal settings. These results support our hypothesis that the parameter-to-compute ratio deemed optimal for bpe models is similarly applicable to {\textsc BLT}, or deviates from it only minimally.

Achieving alignment with bpe scaling patterns necessitates both advances in architecture and the adoption of dynamic patching. As shown in ~\autoref{fig:scaling-compute-opt} (left), we provide a comparison between models utilizing space-patching techniques and Llama 3. Here, we use {\textsc BLT} space-patching as an estimate for SpaceByte \citep{slagle2024spacebyte}, omitting both n-gram embeddings and cross-attention mechanisms. While SpaceByte demonstrates an advantage over Megabyte, its performance still lags considerably behind that of Llama 3. The right panel of \autoref{fig:scaling-compute-opt} depicts the cumulative effect of architectural enhancements together with dynamic patching. On large-scale benchmarks, {\textsc BLT} models match the performance of leading tokenizer-based models such as Llama 3.

Additionally, our findings indicate that for models reliant on tokenization, the selection of tokenizer notably impacts results; specifically, models utilizing the Llama-3 tokenizer achieve higher performance compared to those employing the Llama-2 tokenizer, despite being trained on identical datasets.

In summary, our {\textsc BLT} framework exhibits behavior that falls between Llama 2 and 3 when evaluated with much larger patch sizes. The average token produced by the bpe tokenizers in Llama 2 and 3 contains 3.7 and 4.4 bytes, respectively. By contrast, {\textsc BLT} is capable of displaying comparable scaling patterns while utilizing average patch sizes of 6 and even 8 bytes. Given that inference flop is inversely related to patch size, opting for an 8-byte patch size yields close to a 50\% reduction in inference flops. Moreover, it appears that scaling up patch sizes leads to better model performance as both data and model capacity increase. Although {\textsc BLT} with 8-byte patches initially underperforms relative to bpe Llama 2 at the 1B parameter scale, it ultimately surpasses bpe at the 7B level. This indicates that such large patch sizes could deliver even stronger results at larger scales, suggesting further increases in patch size may be advantageous as training compute and model capacity are expanded.

\subsection{Beyond Compute Optimal Task Evaluations}
\label{section:evals}
To further investigate the effects of scaling, we continue training an 8B {\textsc BLT} model past the conventional compute-optimal threshold using the {\textsc BLT}-1T dataset, which is both larger and higher quality. The resulting model is then evaluated across a diverse set of established benchmarks designed for both classification and generative tasks. We focus specifically on the following types of evaluations to probe common sense reasoning, factual knowledge, and programming competence:
\paragraph{Classification tasks} encompass ARC-Easy (0-shot)~\citep{Eval_ARC}, Arc-Challenge (0-shot)~\citep{Eval_ARC}, HellaSwag (0-shot)~\citep{Eval_hellaswag}, PIQA (0-shot)~\citep{Eval_piqa}, and MMLU (5-shot)~\citep{hendrycks2020measuring}. Our evaluation strategy for these tasks utilizes prompt scoring, where the model’s predicted likelihoods over answer choices are computed and mean accuracy is used as the main metric.
\paragraph{Coding related generation tasks:} We evaluate coding ability using pass@1 results on MBPP (3-shot)~\citep{Eval_MBPP} and HumanEval (0-shot)~\citep{Eval_HumanEval}, thereby gauging how well the language model can produce correct Python code in both few-shot and zero-shot scenarios.

In \autoref{tab:evals}, we present a comparison of three models trained on the {\textsc BLT}-1T dataset: a model utilizing the Llama 3 tokenizer with byte-pair encoding (bpe),\footnote{We select the Llama 3 tokenizer, featuring a 128k vocabulary, due to its superior performance compared to the 32k vocabulary in Llama 2.} alongside two configurations of the {\textsc BLT} model. The first configuration implements space-patching ({\textsc BLT}-Space), while the second adopts an entropy-driven patching mechanism ({\textsc BLT}-Entropy), imposing an approximate monotonicity constraint and resetting the entropy model’s context using line breaks, as detailed in~\autoref{section:entropy-context}. Training is performed under equivalent computational constraints for all models. Notably, for {\textsc BLT}-Entropy, the entropy threshold used during inference is reduced from 0.6 to 0.1, a modification observed to enhance task performance, albeit at the expense of increased inference steps.

The {\textsc BLT}-Entropy model surpasses the Llama 3 model in performance on 4 of the 7 evaluated tasks, despite both models being trained with an equivalent byte count. This superior performance can likely be attributed to (1) more efficient utilization of training computation through dynamic patching, and (2) the explicit modeling of information at the byte level rather than at the token level.

Conversely, {\textsc BLT}-Space lags behind the Llama 3 tokenizer across nearly all benchmarks except one, despite demonstrating a notable decrease in inference flops due to its higher average patch size of 6 bytes. By contrast, both the bpe and entropy-patching approaches exhibit similar average patch sizes of about 4.5 bytes on the same training dataset mix. Given the same training resources, the model employing a larger patch size is able to process 30\% more data than the other two approaches, which could potentially lead {\textsc BLT} to diverge further from the compute-optimal regime.

\subsection{Patches Scale Better Than Tokens} {\textsc BLT} models allow for concurrent scaling of both model capacity and patch dimensions, all while keeping training and inference flops as well as the volume of training data unchanged. This ability to arbitrarily extend patch size distinguishes patch-based approaches, enabling them to sidestep the efficiency limitations that characterize fixed-vocabulary token-based models, as elaborated in Section~\ref{section:bpe-patch}. Utilizing larger patches leads to reductions in computational requirements, freeing resources that can be redirected to expand the global latent transformer, which needs to be invoked less frequently as a result.

A fixed inference scaling study is performed to investigate whether increasing model size, while reducing the number of steps and utilizing larger patches, yields improved performance over smaller models with more steps. Beginning with models containing 400 million and 3.6 billion parameters using the Llama 2 tokenizer, we identify flop-equivalent models employing the Llama 3 tokenizer, as well as {\textsc BLT}-Entropy models that process average patch sizes of 6 and 8 bytes in the training datamix (refer to~\autoref{tab:fixed-inf-params} for additional model specifications). For models configured with a patch size of 8, we employ three encoder layers, in contrast to the single encoder layer used elsewhere. Each model is trained across a range of training flop budgets.

As illustrated in \autoref{fig:fixed_inference_scaling}, {\textsc BLT} models demonstrate superior scaling behavior compared to tokenization-based methods across both inference flop categories. BPE-based models initially yield stronger performance at lower training compute budgets but are soon outperformed by {\textsc BLT} architectures once past the regime of compute-optimality. In practical terms, investing additional resources into the pretraining phase may be advantageous if it results in a more capable model, even when operating under a constant inference budget. A notable case of this approach is the class of 8B models, such as Llama 3.1, which was pretrained on a dataset two orders of magnitude larger than what would be considered optimal given its model size from a compute standpoint.

The point at which {\textsc BLT} surpasses token-based models now lies nearer to the compute-optimal regime in the larger flop class models, shifting from being 3 times to 2.5 times the optimal compute budget. Additionally, within this higher flop regime, the model with patch size 8 exhibits a steeper performance scaling, outperforming the alternatives earlier in the scale-up process. As outlined in Section~\ref{section:parameter-matched-scaling}, increasing patch size yields results more comparable to BPE-based models as model size grows. We partially explain this by noting the reducing proportion of total compute devoted to the byte-level Encoder and Decoder components, which seem to scale more slowly relative to the Latent Transformer. For instance, scaling the overall model parameters by a factor of 20 (from 400M to 8B) results in only an approximate twofold increase in the local parameters of {\textsc BLT}}. This distinction matters because larger patch sizes exclusively impact the flop usage of the patch Latent Transformer, not the byte-level modules. Consequently, this accounts for why the relative size of the {\textsc BLT}-Entropy ps=8 model compared to Llama 2 increased slightly from 1.6x to 1.7x at the larger scale.

To conclude, our investigation into patch-length scaling reveals that the {\textsc BLT} patch-based architecture exhibits superior scaling behaviors when both the patch size and model size are expanded concurrently. These positive trends appear to not only persist but also strengthen as model size increases.

\section{Byte Modeling Improves Robustness} Additionally, we evaluate the robustness of the {\textsc BLT} model relative to token-based architectures that do not process byte-level data directly, and we introduce a method to convert pretrained token-based models to handle bytes. 
\label{sec:robustness}

\subsection{Character-Level Tasks}

One of the original incentives for developing byte-level models was their inherent resilience to noise at the byte level within the input, as well as their capability to recognize the internal structure of tokens—an aspect where models dependent on tokenizers often fall short. To assess these properties, we conduct supplementary experiments using benchmarks designed to test not only the robustness to input noise but also the recognition of characters, spanning English and other languages, and covering elements such as digits and phonemes. The corresponding findings are summarized in Table~\ref{tab:char_tasks}.

\paragraph{Noisy Data} To assess and contrast the robustness of tokenizer-based models with {\textsc BLT}, we construct altered versions of the benchmark classification tasks introduced in Section~\ref{section:evals}. We utilize five separate character-level noise injection methods to manipulate the original text: (a) \textit{AntSpeak}: This approach changes all letters to uppercase and places a space between each character. (b) \textit{Drop}: Eliminates 10\% of the text's characters at random. (c) \textit{RandomCase}: Randomly switches 50\% of the letters to uppercase and the remaining 50\% to lowercase throughout the entire sample. (d) \textit{Repeat}: Selects 20\% of the characters and repeats each up to four times. (e) \textit{UpperCase}: Converts the entire string to uppercase. Each noising scheme is independently applied to either the prompt, the completion, or both components during evaluation, and mean performance scores are reported. Table~\ref{tab:char_tasks} presents outcomes on the perturbed HellaSwag~\citep{Eval_hellaswag} dataset, demonstrating that {\textsc BLT} consistently surpasses tokenizer-based counterparts in robustness. On average, it achieves an 8-point lead over the model trained on identical data and even shows gains relative to Llama 3.1, which is trained on a far larger corpus.

\paragraph{Phonology - Grapheme-to-Phoneme (G2P)} We evaluate the effectiveness of {\textsc BLT} in converting sequences of graphemes (the written symbols constituting a word) into their corresponding phonemic representations (the sounds of the word). Table \ref{tab:char_tasks} displays the performance of the models on the G2P task using a 5-shot configuration with Phonology Bench~\citep{suvarna2024phonologybench}. Our findings indicate that {\textsc BLT} achieves superior results compared to the baseline model that uses the Llama 3 1T tokenizer for this task.

\paragraph{CUTE}
To evaluate the character-level capabilities of {\textsc BLT}, we conduct experiments on the CUTE benchmark~\citep{edman2024cute}, which contains tasks organized into three main categories: compositional understanding, orthographic similarity comprehension, and sequence manipulation. This suite of tasks presents a notable difficulty for most models reliant on tokenization, as although these models seem aware of the spelling of their constituent tokens, they often fail to harness this information effectively in text manipulation. Results in Table~\ref{tab:char_tasks} reveal that {\textsc BLT}-Entropy surpasses both BPE Llama 3 baselines by over 25 points on the CUTE benchmark. Notably, our approach exhibits near-perfect performance in tasks requiring character manipulation, attaining 99.9\% accuracy on both spelling-related tasks. These substantial improvements—despite {\textsc BLT} using only 1/16th the data of Llama 3.1—suggest that acquiring character-level knowledge remains particularly difficult for BPE-based architectures. Figure \ref{fig:cute_examples} showcases examples where the Llama 3 tokenizer-based model underperforms in contrast to the strong results achieved by our {\textsc BLT} model. The only tasks in which BPE-based models demonstrate an advantage are word deletion and insertion. While such manipulations may be less direct for a byte-level system, the performance gap is modest, and constructing word-level representations from characters may prove less challenging than the reverse. Our experiments utilize consistent evaluation procedures across all tasks, adopting the original Huggingface prompts. Further prompt engineering might provide additional gains for the BPE models.

\paragraph{Low Resource Machine Translation} We assess the performance of {\textsc BLT} for translation tasks both to and from twenty-one lower-resource languages, representing six widely recognized language families and a range of scripts, as presented in the FLORES-101 benchmark~\citep{goyal2022flores}. Detailed SentencePiece BLEU scores are summarized in Table~\ref{tab:flores}. The findings indicate that {\textsc BLT} achieves better results than a counterpart model utilizing the Llama 3 tokenizer, outperforming it by 2 BLEU points when translating into English and by 0.5 BLEU points for translations from English. In high-resource language pairs, {\textsc BLT} matches or marginally exceeds Llama 3 in performance. Notably, for several language pairs situated in lower-resource families, {\textsc BLT} consistently surpasses Llama 3, highlighting the strengths of byte-level modeling in handling less common byte sequences and generalizing to underrepresented languages.

\subsection{Training {\textsc BLT} Using Llama 3 Initialization}
\label{sec:distillation}
We investigate an approach where {\textsc BLT} models benefit from the parameters of pre-existing tokenizer-based models, thereby enabling more efficient and accelerated convergence during training. To accomplish this, we initialize the global transformer parameters within {\textsc BLT} using those from a pre-trained Llama 3.1 model. The training procedure then involves updating the global transformer weights at a learning rate set to one-tenth of that applied to the local encoder and local decoder, for the optimal number of training steps established for Llama 3. We summarize a comparison against a baseline {\textsc BLT} model in Table~\ref{tab:distill}. The results indicate a marked advantage of the {\textsc BLT} model initialized from Llama 3.1, which exceeds the performance of both Llama 3 and the baseline {\textsc BLT}—all models receiving equal compute resources. Furthermore, even relative to our {\textsc BLT}-Entropy model (results in Table~\ref{tab:evals}), which underwent training on a substantially larger 1T-token corpus, the Llama 3.1-initialized {\textsc BLT} achieves better scores on the MMLU benchmark. This demonstrates that this transfer approach can dramatically reduce the computational requirements for training while maintaining or improving model quality.

This approach can alternatively be interpreted as transforming models that rely on tokenizers into models that do not, thereby repurposing a pre-trained LLaMA 3.1 model as a {\textsc BLT} model. To ensure a thorough evaluation, we report results for the original LLaMA 3.1 trained on 15T tokens in Table \ref{tab:distill} and compare its performance to that of the {\textsc BLT} model derived from LLaMA 3. While our method shows only minor decreases in performance on MMLU and HumanEval, it experiences more pronounced reductions in other evaluation tasks. This indicates that there is room for improvement in fully utilizing the benefits of pre-trained models—particularly through refining data composition and adjusting hyperparameters.

\section{Ablations and Discussion}
\label{section:ablations}
Here, we examine ablation studies that support the design decisions made for {\textsc BLT}, as well as evaluate the patching methodology and selected hyper-parameters used in training the 8B parameter {\textsc BLT} model on the {\textsc BLT}-1T dataset.

\paragraph{Entropy Model Hyper-parameters} In order to investigate how adjustments to entropy model size and context window length influence scaling behavior, we train several byte-level entropy transformer models, varying their parameters between 1m and 100m and selecting context window lengths in the range of 64 to 512. For this analysis, we display bits-per-byte (bpb) against the number of training FLOPs by generating scaling law curves, leveraging our $400m$ and $1b$ {\textsc BLT} models, both of which are trained on the Llama-2 dataset, as illustrated in \autoref{fig:entropy_ablation}. The results show a direct relationship between scaling outcomes and increases in both model size and context window length, although the rate of improvement notably slows for model sizes exceeding 50m parameters.

\paragraph{Types of Patching}

We conduct an ablation study on the four patching approaches described in Section~\ref{section:patching}: (1) Strided Patching utilizing strides of 4 and 6, (2) Whitespace-based Patching, (3) Patching informed by the BPE Tokenizer of Llama 3, and (4) Entropy-based patching with the aid of a compact byte-level language model.

Although dynamic patching effectively shortens the sequence length, we regulate the sequence length across all patching methods to ensure that the context window remains consistent. This approach ensures that, on average, each model is exposed to an identical quantity of bytes per sequence during both training and inference, thereby eliminating potential confounds related to differing context sizes. The outcomes of these ablation studies are illustrated in Figure~\ref{fig:scaling-compute-opt}. Notably, every alternative patching scheme surpasses static patching in performance, and among these, space patching closely matches the effectiveness of dynamic entropy based patching.

In \autoref{tab:evals_ablation}, we report benchmark results for {\textsc BLT} models, examining the performance of tokenizer-based models alongside those trained using space patching and entropy-based patching strategies. All models were trained on the Llama 2 dataset for the optimal number of training steps, as described in~\citep{dubey2024llama}. While space patching offers a more straightforward approach that bypasses the need for real-time entropy modeling during training, our findings reveal that the benefits of entropy-based patching observed in scaling behavior (see Section~\ref{section:scaling}) also extend to performance on downstream benchmarks.\footnote{Space patching results stem from prior experiments without cross-attention, but comparable outcomes are seen with the inclusion of cross-attention.}

\paragraph{Cross-Attention} 
\autoref{tab:crossattn_ablations_1b} presents an ablation study exploring the impact of introducing cross-attention at different locations within the encoder and decoder modules of {\textsc BLT}. For the encoder side, cross-attention mechanisms are assessed using three distinct query initialization strategies: (1) a uniform learned embedding applied across all global states, (2) hash-based embeddings derived from the bytes present in the patch, and (3) representations obtained by pooling the encoder's hidden states corresponding to the patch bytes at each respective encoder layer.

Our results indicate that incorporating cross-attention within the \textit{decoder} yields the highest effectiveness. In contrast, integrating cross-attention into the encoder brings only marginal gains, and these appear primarily when queries are initialized via pooling. Furthermore, our analysis shows that cross-attention provides notable benefits on Common-Crawl datasets, with the advantage being more pronounced for larger patch sizes.

\paragraph{n-gram Hash Embeddings} 

We experiment with various n-gram hash embedding vocabulary sizes—specifically 0, 100K, 200K, and 400K—and report these findings in Table~\ref{tab:ngram_ablations_1b}. Our analysis reveals that the introduction of hash embeddings consistently benefits performance across all evaluated domains, with notable improvements on Wikipedia and Github (a 0.04 bpb gain, versus 0.01 bpb after 15k training steps at the 8B parameter scale). At the same 8B scale, reducing the number of hashes from 500K to 300K yields only a minimal performance reduction of roughly 0.001 bpb after 15k steps. These outcomes suggest that hash embeddings are crucial for enabling {\textsc BLT} models to achieve parity with traditional tokenizer-based approaches. Nonetheless, the performance enhancement plateaus beyond 300K hashes, signaling diminishing returns. Furthermore, the results suggest the improvements from hash embeddings and cross-attention mechanisms are largely independent, as they benefit performance across distinct datasets.

\paragraph{Local Model Hyperparameters} Table \ref{tab:local_layers} presents an ablation study exploring different configurations for the number of layers used in the local encoder and decoder. When utilizing hash n-gram embeddings, {\textsc BLT} achieves strong performance even when the encoder is minimal—specifically, comprising only a single layer—while the decoder is comparatively deeper.

\section{Related Work}
\label{section:related_work}

\textbf{Character-Level RNNs:} The task of Character Language Modeling has long been prominent in neural approaches~\citep{sutskever2011generating,mikolov2012subword,graves2013generating}, largely because it naturally accommodates the modeling of out-of-vocabulary words, thereby removing the need for fallback strategies. \cite{kim2016character} introduced a model that operates solely on character inputs, employing convolutional layers and highway networks as input encoders into an LSTM-based RNN architecture; their approach achieved parity with the then state-of-the-art RNN language models in English and surpassed them for languages with complex morphology—a recognized strength of character-level large language models. Building on this, \cite{kenter2018byte} applied byte-level LSTMs for machine comprehension tasks, reporting superior results over word-level systems, particularly in morphologically rich languages such as Turkish and Russian. In related work, \cite{zhang2015character} adopted character-level convolutional models for classification, demonstrating advantages over word-based models in specific tasks. Expanding on these ideas, \citet{chung2019hierarchical} proposed hierarchical LSTM architectures, introducing boundary detectors at each stage to automatically infer textual hierarchies, further boosting performance for character-level language modeling. Alternatively, \cite{kalchbrenner2016neural} presented ByteNet, which replaced attention mechanisms with convolutional networks at the character level for machine translation tasks.

\textbf{Character-Level Transformers:} The introduction of attention mechanisms in transformer architectures~\citep{vaswani2017attention} and the adoption of subword tokenization~\citep{sennrich-etal-2016-neural} brought about notable gains in neural language modeling and evaluation benchmarks. Nonetheless, employing word or subword representations imposes an inherent inductive bias regarding the abstraction level at which the models function. To unify the advantages of transformers with encouraging early work in character-level language modeling, \citet{al2019character} employed very deep transformer networks and leveraged auxiliary loss functions, successfully training transformer models that surpassed the prior LSTM-based character-level language models. Yet, these models continued to lag considerably behind their word-level LLM counterparts. Similarly, GPT-2~\citep{radfordgpt2} reported that for large-scale datasets such as the 1 Billion Word Benchmark, byte-level language models fell short of matching the performance achieved by word-level language models.

\cite{Choe2019BridgingTG} provided evidence that transformer-based large language models (LLMs) operating at the byte level can surpass subword-based LLMs when parameter counts are similar; however, this advantage comes at the cost of significantly increased computational requirements and extended training durations. In a related vein, \cite{el2020characterbert} introduced CharFormer, a BERT-based architecture that generates word-level representations by applying convolutional layers to character embeddings. Although this approach yields notable gains within the medical domain, it also results in considerable computational overhead. Likewise, \cite{clark2022canine} proposed CANINE, a 150M-parameter model designed to encode character sequences directly. At the model's core lies a deep transformer architecture reminiscent of the global module presented in our work, paired with a localized transformer mechanism and strided convolutions for effective input downsampling. CANINE achieved superior performance compared to mBERT, the token-level encoder-only counterpart, across downstream multilingual benchmarks. The ByT5 model~\citep{xue2022byt5} further investigated strategies for byte-level encoder-decoder architectures, intentionally foregoing any form of patch-based processing. Although ByT5 demonstrated both heightened resilience to input noise and strong performance relative to tokenization-based models using only a quarter of the usual data, its absence of patching necessitated resource-intensive attention computations over each byte, dramatically increasing computational load. Directly processing bytes rather than subword representations substantially raises input sequence length, posing challenges for the scalable and efficient training of byte-level models. Recent developments utilizing the Mamba Architecture~\citep{gu2023mamba}, characterized by the ability to preserve a fixed-size memory state across long contexts, have enabled \cite{wang2024mambabyte} to train a patch-free, byte-level Mamba model. Their work demonstrated superior results in terms of bits-per-byte compared to byte-level transformer baselines, specifically at the 350M parameter scale in flop-controlled scenarios and across multiple datasets.

\textbf{Patching-based approaches:} Leveraging patching methods has the potential to significantly reduce the high computational cost in terms of flops typically associated with byte-level LLMs, all while possibly maintaining their effectiveness. A number of studies have reported promising results with smaller models and limited training data. For instance, \cite{nawrot-etal-2022-hierarchical} investigate static patching approaches utilizing downsampling and upsampling operations, and introduce the hourglass transformer, which surpasses other byte-level baselines at the 150M parameter scale. Building on this, \cite{nawrot-etal-2023-efficient} enhance the method by introducing dynamic patching mechanisms, such as an end-to-end trained boundary predictor, a boundary predictor guided by certain tokenization methods, and an entropy-driven patching strategy reminiscent of {\textsc BLT}. Their results demonstrate that these techniques can outperform standard transformers of the era on language modeling benchmarks with a 40M parameter model trained on approximately 400M tokens. In another line of work, \cite{lester2024training} explore training LLMs on sequences compressed via arithmetic coding to surpass the compression rates achievable by BPE, and implement an equal-info windows method. While this allows them to achieve better results than byte-level models for language modeling, their approach still falls short of models trained on subword representations.

This study is principally informed by and closely aligns with the MegaByte~\citep{yu2023megabyte} architecture, a decoder-only, causal large language model that employs a fixed, static strategy for patching by concatenating representations to segment bytes into patches. It subsequently utilizes a local decoder model to revert patches back to their byte forms. The authors of MegaByte show that the model attains performance competitive with tokenizer-based approaches when evaluated at a 1B parameter scale on datasets containing roughly 400B bytes. In our analyses, we directly compare against MegaByte and observe that its static patching mechanism falls short of the leading compute-efficient, optimally trained tokenizer-based models under flop constraints; through our investigations, we illustrate how {\textsc BLT} addresses this shortcoming. Furthermore, \cite{slagle2024spacebyte} similarly note MegaByte's limitations and propose advancing the static patching framework by incorporating patching on whitespace or space-like bytes, in addition to integrating a local encoder. Their findings suggest that this extension provides improvements in compute-controlled scenarios for certain data sources (e.g., Github and arXiv) at the 1B parameter scale. Our own experiments with this configuration reinforce the view that while these modifications offer benefits, significant architectural advancements are still required for byte-level models to scale and perform on par with the latest token-based state-of-the-art systems such as Llama 3.

\section{Limitations and Future Work}

In this study, we base our architectural decisions on models trained for the optimal number of steps, following the approach established for Llama~3~\citep{dubey2024llama}. Nevertheless, it is important to note that these scaling laws were originally derived for BPE-level transformers and may not yield the most efficient ratios of data to parameter sizes when applied to {\textsc BLT}. Determining precise scaling laws for {\textsc BLT} remains an open question for future investigation and could reveal even more advantageous scaling patterns suited to our proposed architecture. Furthermore, since much of our experimentation has been limited to model sizes up to 1B parameters, it is conceivable that ideal architectural configurations may shift as we expand to models with 8B parameters or larger, potentially leading to enhanced performance at greater scales.

Current transformer libraries and codebases primarily optimize for tokenizer-based transformer models. Although we conduct theoretically flop-matched experiments and leverage specialized efficient methods, like FlexAttention, to support non-standard transformer layers, our current implementations might still lag behind tokenizer-based models in actual wall-clock performance and could be improved further with additional optimizations.

Whereas {\textsc BLT} employs an independently trained entropy model for patching, an intriguing avenue for future research involves integrating the patching model into a fully end-to-end training pipeline. In Section \ref{sec:distillation}, we describe preliminary experiments that suggest promising outcomes for the process of ``byte-ifying'' tokenizer-based architectures like Llama 3, which have been exposed to over 10T tokens, by initializing and subsequently freezing the weights of their global transformer components. Expanding on this line of investigation may yield approaches that both preserve the advantages conferred by bytefication and potentially improve upon the performance of tokenizer-based models, all without the need for retraining from the ground up.

\section{Conclusion}
\label{section:conclusion}

This work introduces the Byte Latent Transformer (\textbf{BLT}), a novel model architecture that challenges the prevailing reliance on fixed-vocabulary tokenization methods in large language models. By employing a flexible, trainable mechanism to group bytes into adaptive patches, {\textsc BLT} dynamically allocates computational effort according to the intricacy of the input, thereby enhancing both efficiency and model robustness. Through a comprehensive scaling analysis, we show that {\textsc BLT} matches the capabilities of tokenization-dependent architectures such as Llama 3 for configurations up to 8B parameters and 4T bytes, and can achieve up to a 50\% reduction in inference flops by tolerating modest degradations in standard evaluation scores. Additionally, {\textsc BLT} enables an unexplored scaling direction: simultaneously enlarging both the model and patch sizes while keeping inference costs constant, which proves beneficial in many real-world computational environments. By processing sequences at the byte level, {\textsc BLT} enhances robustness to noisy data and exhibits improved modeling of rare and complex sub-word patterns. Collectively, these findings identify {\textsc BLT} as a compelling alternative to conventional tokenization-based techniques, delivering a versatile, scalable, and robust approach for advancing efficient language modeling.

\end{document}